\documentclass[11pt]{article}
\usepackage[margin=25mm]{geometry}
\usepackage{fontspec,unicode-math}
\setmainfont{Latin Modern Roman}
\setmathfont{Latin Modern Math}
\usepackage{amsmath,amssymb,mathtools}
\usepackage{xurl,hyperref}
\hypersetup{colorlinks=true,urlcolor=blue,linkcolor=blue}
\setcounter{secnumdepth}{0}
\setlength{\parindent}{0pt}
\setlength{\parskip}{5pt}
\setlength{\emergencystretch}{3em}
\newcommand{\sref}[2]{\hyperlink{source:#1:#2}{[S#2]}}
\newcommand{\cataloguescope}[1]{\paragraph{Source question.}#1}
\newcommand{\floor}[1]{\left\lfloor#1\right\rfloor}
\begin{document}
\section{Minimum Crossing Realizations of Spatial Polygon Angle Sequences}\label{statement}

\subsection{Definitions and mathematical statement}

\paragraph{Shared notation.}Standard number systems are denoted by $\mathbb N,\mathbb Z,\mathbb Q,\mathbb R,\mathbb C$; index ranges, computational models and specialized notation follow the cited problem source and the definitions below.


A spatial polygon is a closed polygonal chain in $\mathbb R^3$ with positive edge lengths. Its unsigned turning angles are the angles between consecutive oriented edge directions, in $(0,\pi)$. An angle sequence is realizable if such a polygon exists. In the cited work, crossing counts refer to self-intersections of the polygonal chain, with generic realizations excluding coincident vertices, overlaps, and degenerate intersections; they are not crossings of an arbitrary planar projection.
\paragraph{Statement recorded in the supplied source.}
Given a realizable angle sequence for a polygon in $\mathbb{R}^3$, can an efficient algorithm compute the minimum number of crossings among its generic realizations? \sref{statement}{2}
\paragraph{Scope qualification.}
The relevant convention is the unsigned turning-angle convention of the cited polygon paper. ``Crossings'' here does not ask for the knot-theoretic crossing number. The source does not prescribe a bit model for real-valued angles.
\paragraph{Source status.}
The supplied catalogue labels this record open. Its source-status metadata was checked against the supplied snapshot on 2026-10-11. This import does not independently establish current open status or mathematical correctness. The catalogue literature assessment is dated 2026-08-22.
\subsection{Short English statement}
Find a spatial polygon with the prescribed angles and as few self-intersections as possible.
\subsection{Sources}
\begin{itemize}
\item[\hypertarget{source:statement:1}{S1}] Ulam AI and the UnsolvedMath Contributors, UnsolvedMath: A Curated Collection of Open Mathematics Problems, catalogue entry. Catalogue curation and metadata: CC BY 4.0. The supplied source question is retained; definitions and the English summary are editorial additions.
\url{https://huggingface.co/datasets/ulamai/UnsolvedMath}.
\item[\hypertarget{source:statement:2}{S2}] Discrete Geometry (2021), Oberwolfach Reports, DOI 10.4171/owr/2020/30, pp. 1469--1528
\url{https://doi.org/10.4171/owr/2020/30}.
\item[\hypertarget{source:statement:3}{S3}] A. Efrat, R. Fulek, S. Kobourov and C. D. Tóth, Polygons with Prescribed Angles in 2D and 3D, arXiv:2008.10192
\url{https://arxiv.org/abs/2008.10192}.
\textit{Underlying Oberwolfach Reports material retains its source-specific licence, including CC BY-SA 4.0 where applicable; see THIRD\_PARTY\_NOTICES.md.}
\end{itemize}
\paragraph{Preparation and review.}
Original question attributed to the cited source and UnsolvedMath contributors. Codex (GPT-6; exact serving revision not exposed) assisted with definitions, English summaries, source inspection and formatting. No new human mathematical authorship or independent proof verification is asserted. Catalogue identity checked against the complete supplied registry and existing repository statements; known flagged repetitions and obvious identical conjectures excluded. Full mathematical-equivalence/literature review has not been performed. Statement preparation and final review followed contribution prompts C and D; automated checks verify structure and identity, not open status.
\label{end:statement}
\end{document}
